\documentclass[10pt,a4paper]{amsart}
\usepackage[margin=19mm]{geometry}
\usepackage{amsmath,amssymb,mathtools}
\usepackage[scr=rsfs,cal=euler]{mathalfa}
\usepackage{xfrac}
\usepackage[unicode,hidelinks,bookmarks=false]{hyperref}
\newcommand{\ie}{i.e.\ }
\newcommand{\cf}{cf.\ }
\newcommand{\resp}{respectively\ }
\newcommand{\Subsection}{\subsection}
\providecommand{\nobreakdash}{\nobreak\hbox{-}}
\setlength{\emergencystretch}{2em}
\multlinegap0pt
\usepackage[shortlabels]{enumitem}
\usepackage{amssymb}
\usepackage{mathtools}
\usepackage{xcolor}
\usepackage{algorithm}\usepackage{algpseudocode}
\newtheorem{assumption}{Assumption}
\newtheorem{definition}{Definition}
\newtheorem{theorem}{Theorem}
\title{Nonsmooth Obstacles and Killed Resolvents in Reflected Stochastic Control}
\author{Louis Shuo Wang, Jiguang Yu}
\date{Source: arXiv:2606.21575v2 (CC BY 4.0)}
\begin{document}
\maketitle

\noindent\emph{Source and licence.} Nonsmooth Obstacles and Killed Resolvents in Reflected Stochastic Control. Version arXiv:2606.21575v2 (CC BY 4.0), distributed by arXiv under the Creative Commons Attribution 4.0 International licence, http://creativecommons.org/licenses/by/4.0/, as recorded in the arXiv licence metadata for this exact version and shown on the abstract page https://arxiv.org/abs/2606.21575v2. Full licence notice: \texttt{licenses/arxiv-2606.21575-CC-BY-4.0.txt}.\par\smallskip

\noindent\textit{Editorial scope.} Counted unit: the article's theorem thm:verification-boundary (source0.tex lines 614--632). The article's complete original proof of it is delivered with it (source0.tex lines 634--673, one \texttt{proof} environment of 40 lines). No mathematics is written, repaired or re-derived: the statement and the proof are the article's own lines, in the article's own notation, and no retained line is substituted or re-flowed.  Retained uncounted as the source-local context the proof needs: lines 262--268; lines 585--587; lines 598--609.  Nothing was omitted from any retained span, so the package carries no omission record.  No retained line contains a citation, so the wrapper carries no bibliography and no bibliography entry.  No retained line contains a figure environment, an image-inclusion command or drawing code, so no image is delivered and the package declares no asset dependency.  Editorial formatting only, in addition: an AMS \texttt{amsart} preamble that restates the article's own declarations and macros used by the retained lines, namely the environments \texttt{assumption}, \texttt{definition}, \texttt{theorem} on separate counters, together with the standard packages those lines need.  The retained environments print in this document as Assumption 1, Definition 1, Theorem 1 in order of appearance, whereas the article prints the same items with the numbering of its own full document.  The complete native source of the article as distributed in the arXiv e-print for this exact version is bundled unchanged as \texttt{sources/203278/source0.tex}.
\begin{assumption}[Lyapunov growth control] \label{ass:lyapunov} 
There exist a function \(\Psi\in C^2(\mathbb R_+^2)\), constants \(K_\Psi,\lambda_\Psi>0\), and an exponent \(m\ge1\), such that 
\[ 1+|x|^m\le K_\Psi \Psi(x), \qquad G(x)+c(x)\le K_\Psi \Psi(x), \] 
and 
\[ \mathcal L\Psi(x)\le \lambda_\Psi \Psi(x)+K_\Psi, \qquad x\in(0,\infty)^2 . \] 
Moreover, \(r>\lambda_\Psi\). 
\end{assumption}

\begin{equation} \label{eq:Ub-definition} 
U_b(x):=G(x)-R_r^{\mathcal C_b}\Gamma(x). 
\end{equation} 

\begin{definition}[Admissible candidate value] \label{def:admissible-candidate-value} 
Let \(b\) be locally Lipschitz and define \(U_b\) by \eqref{eq:Ub-definition}. We say that \(U_b\) is an admissible candidate value if the following hold: 
\begin{enumerate}[label=(\roman*)] 
\item \(U_b\) is continuous on \(\mathbb R_+^2\); 
\item \(U_b\) belongs locally to the generalized It\^{o} class on \(\mathbb R_+^2\); 
\item \(U_b\) has growth controlled by the Lyapunov function of Assumption~\ref{ass:lyapunov}; 
\item the reflected Neumann condition holds on the coordinate faces in the trace sense: 
\[ 
\partial_i U_b=0 \qquad\text{on }\{x_i=0\},\quad i=1,2; 
\] 
\item the stopped process \[ \left\{ e^{-r(t\wedge\tau_b)}U_b(X_{t\wedge\tau_b}) - \int_0^{t\wedge\tau_b}e^{-rs}c(X_s)\,ds \right\}_{t\ge0} \] is of class \(D\) after localization, so that optional sampling is valid. \end{enumerate} 
\end{definition} 

We now state the main verification theorem. It converts a candidate boundary into a verified optimizer. \begin{theorem}[Verification of a candidate boundary] \label{thm:verification-boundary} Let \(b\) be a locally Lipschitz boundary and define 
\[ \mathcal D_b=\{(x_1,x_2):x_2\ge b(x_1)\}, \qquad \mathcal C_b=\mathbb R_+^2\setminus\mathcal D_b. \] 
Set 
$\displaystyle U_b=G-R_r^{\mathcal C_b}\Gamma$.  
Assume that: 
\begin{enumerate}[label=(\roman*)]
\item \(U_b\ge G\) on \(\mathbb R_+^2\); 
\item \(U_b=G\) on \(\mathcal D_b\); 
\item \(U_b>G\) on \(\mathcal C_b\); 
\item \(U_b\) satisfies the reflected Neumann condition on the coordinate axes; 
\item \((r-\mathcal L)U_b+c\ge0\) in the sense of measures on \(\mathbb R_+^2\); 
\item \(U_b\) has admissible growth and belongs to the generalized It\^{o} class in the sense of Definition~\ref{def:admissible-candidate-value}. 
\end{enumerate} 
Then 
$U_b(x)=V(x)$ for $x\in\mathbb R_+^2$, 
and 
$\tau_b=\inf\{t\ge0:X_t\in\mathcal D_b\}$ 
is an optimal stopping time. 
\end{theorem} 

\begin{proof} 
We split the proof into the usual super-solution and attainability steps. First, we prove that \(U_b\) dominates the value. Since 
$(r-\mathcal L)U_b+c\ge0$ in the sense of measures, the generalized It\^{o} formula gives, after localization, 
\[
e^{-r(t\wedge\tau)}U_b(X_{t\wedge\tau}) = U_b(x) + \int_0^{t\wedge\tau} e^{-rs} \big(\mathcal L U_b-rU_b\big)(X_s)\,ds + M_{t\wedge\tau} + A_{t\wedge\tau}^{\rm sing},
\]
where \(M\) is a local martingale and \(A^{\rm sing}\) denotes the measure-valued contribution. The measure inequality 
$\displaystyle (r-\mathcal L)U_b+c\ge0$ 
is equivalent to  
$\displaystyle \mathcal L U_b-rU_b\le c $
in the weak sense. The reflected Neumann condition eliminates the boundary local-time terms on the coordinate axes. Hence 
$\displaystyle e^{-r(t\wedge\tau)}U_b(X_{t\wedge\tau}) - \int_0^{t\wedge\tau}e^{-rs}c(X_s)\,ds $
is a supermartingale. By optional sampling and the growth condition, for any admissible stopping time \(\tau\), 
\[ 
U_b(x) \ge \mathbb E_x \left[ e^{-r\tau}U_b(X_\tau) - \int_0^\tau e^{-rs}c(X_s)\,ds \right]. 
\] 
Using \(U_b\ge G\), we obtain 
\[ 
U_b(x) \ge \mathbb E_x \left[ e^{-r\tau}G(X_\tau) - \int_0^\tau e^{-rs}c(X_s)\,ds \right]. 
\] 
Taking the supremum over all stopping times gives 
$U_b(x)\ge V(x)$. 
Second, we prove the reverse inequality. Let 
$\tau_b=\inf\{t\ge0:X_t\in\mathcal D_b\}$. 
On \(\mathcal C_b\), the definition 
$U_b=G-R_r^{\mathcal C_b}\Gamma$ 
and the killed-resolvent construction imply 
$\mathcal L U_b-rU_b-c=0 $
in the continuation region in the weak sense. Therefore, the process 
$\displaystyle e^{-r(t\wedge\tau_b)}U_b(X_{t\wedge\tau_b}) - \int_0^{t\wedge\tau_b}e^{-rs}c(X_s)\,ds$ 
is a martingale after localization. Letting \(t\to\infty\), using the admissible growth condition and the class \(D\) property, yields 
\[ 
U_b(x) = \mathbb E_x \left[ e^{-r\tau_b}U_b(X_{\tau_b}) - \int_0^{\tau_b}e^{-rs}c(X_s)\,ds \right]. 
\] 
Since \(X_{\tau_b}\in\mathcal D_b\) and \(U_b=G\) on \(\mathcal D_b\), we get 
\[ 
U_b(x) = \mathbb E_x \left[ e^{-r\tau_b}G(X_{\tau_b}) - \int_0^{\tau_b}e^{-rs}c(X_s)\,ds \right] \le V(x). 
\] 
Combining the two inequalities gives \(U_b=V\), and the displayed equality shows that \(\tau_b\) is optimal. 
\end{proof} 

\end{document}
